\begin{abstract}
We present a reinforcement-learning framework for training small language
models to reason on verifiable problems. The agent performs approximate
$Q$-learning over an LLM-induced context Markov decision process: states are
structured contexts, actions are a small set of coarse cognitive verbs the
model realizes as text, transitions are sampled from the model's responses,
and reward is anchored by a non-LLM verifier (a symbolic solver, unit-test
harness, or proof checker) rather than a learned critic. The action-value
function $\Qpsi(s, a)$ is a multi-head $Q$-network over the hidden state of a
frozen base model. Monte Carlo Tree Search enters not as the core learning
algorithm but as a policy-improvement operator: it produces locally improved
estimates $\Qmcts(s, a)$ and visit-count policies $\pimcts$ that serve as
stronger training targets for expert iteration, and it doubles as a
test-time-compute knob at inference. Keeping reward verifier-anchored rather
than learned reduces the Goodhart attack surface. We give the framework, the
architecture, and a pre-registered experimental plan with falsifiable
hypotheses (a fixed baseline grid, a FLOP-accounted compute axis that charges
every LLM forward inside the search, and a Holm-Bonferroni-corrected criterion
committed before training). We do not report benchmark numbers: this is the
plan, not results. A preliminary feasibility probe established that the loop is
implementable end-to-end and surfaced the methodological precautions the plan
inherits (train/held-out discipline, per-class composition tracking, a
pairwise-recommended-verb target-stability diagnostic, and honest per-decision
compute accounting with an observable, resumable run harness); that probe
predates the current action space and its specific numbers are superseded. The
large-scale benchmark claim is deferred to the pre-registered study.
\end{abstract}
